%%%%%%%%%%%%%%%%%%%%%%%%%%%%%%%%%%%%%%%%%%%%%%%%%%%%%%%%%%%%%%%%%%%%%%%%
% Plantilla TFG/TFM
% Universidad de Murcia. Facultad de Informática
%%%%%%%%%%%%%%%%%%%%%%%%%%%%%%%%%%%%%%%%%%%%%%%%%%%%%%%%%%%%%%%%%%%%%%%%

\begin{otherlanguage}{english}
\addchap{Extended Abstract}

% El extended abstract se escribe en inglés. Consulta en la normativa
% vigente de la Facultad la extensión exigida (la plantilla original
% indicaba 2000 palabras).

This extended abstract must explain, clearly and concisely, what the project is about, why this topic was chosen and what its main objective is. It must describe the main methods, tools or approaches used during the development, highlight the most relevant results and summarise the main conclusions. It may also discuss the importance of the work and possible applications or future lines of development.

It should be written so that any reader can understand the content and scope of the work without reading the full document.
\end{otherlanguage}
